% =====================================================================
%  THÈME 1 — Pondération des équations chimiques
%  Niveau MOYEN — Exercices
% =====================================================================
\documentclass[11pt,a4paper]{article}
\usepackage[utf8]{inputenc}
\usepackage[T1]{fontenc}
\usepackage{lmodern}
\usepackage[french]{babel}
\usepackage{amsmath,amssymb,mathtools}
\usepackage{icomma}
\usepackage[version=4]{mhchem}
\usepackage{siunitx}
\sisetup{output-decimal-marker={,},per-mode=symbol}
\usepackage{xcolor}
\usepackage{tikz}
\usepackage[most]{tcolorbox}
\usepackage{enumitem}
\usepackage{array}
\usepackage{fancyhdr}
\usepackage[a4paper,margin=2.1cm,headheight=26pt,headsep=8pt]{geometry}

\definecolor{primary}{RGB}{142,93,168}
\definecolor{primarydark}{RGB}{82,45,108}
\definecolor{moycol}{RGB}{198,120,9}

\pagestyle{fancy}\fancyhf{}
\fancyhead[L]{\small\textcolor{primarydark}{\textbf{Fiche 5} — Thème 1 : Pondération}}
\fancyhead[R]{\small\textcolor{moycol}{\textsc{Moyen}}}
\fancyfoot[C]{\small\thepage}
\renewcommand{\headrulewidth}{0.8pt}
\renewcommand{\headrule}{\hbox to\headwidth{\color{primary}\leaders\hrule height \headrulewidth\hfill}}

\newtcolorbox[auto counter]{exo}[1][]{enhanced,breakable,boxrule=1pt,arc=3pt,
  left=3mm,right=3mm,top=2.4mm,bottom=2.4mm,colback=moycol!6,colframe=moycol,
  fonttitle=\bfseries,coltitle=white,
  title={Exercice~\thetcbcounter\ \ifx\relax#1\relax\relax\else\textmd{--- \itshape #1}\fi}}

\setlength{\parindent}{0pt}
\setlength{\parskip}{3pt}

\begin{document}

\begin{center}
{\LARGE\bfseries\color{primarydark}Thème 1 — Pondération}\\[1pt]
{\large\color{moycol}\textbf{Niveau Moyen}}\\[2pt]
{\itshape Objectif : gérer les coefficients fractionnaires, les groupements polyatomiques et les charges.}\\
{\color{primary}\rule{0.6\linewidth}{1pt}}
\end{center}

\begin{exo}[combustion du propane]
Le propane (gaz de camping) brûle selon
\[ \ce{C3H8 + O2 -> CO2 + H2O}\qquad\text{(non pondérée).} \]
Pondère l'équation (carbone, puis hydrogène, puis oxygène) et donne la somme des coefficients.
\end{exo}

\begin{exo}[combustion de l'éthane : gérer une fraction]
L'éthane brûle selon
\[ \ce{C2H6 + O2 -> CO2 + H2O}\qquad\text{(non pondérée).} \]
\textbf{a)} Équilibre d'abord le carbone et l'hydrogène. Quel coefficient \emph{fractionnaire}
obtiens-tu alors devant \ce{O2} ?\\
\textbf{b)} Multiplie toute l'équation pour revenir à des entiers, puis donne la somme des coefficients.
\end{exo}

\begin{exo}[un groupement polyatomique intact]
La réaction de l'hydroxyde d'aluminium avec l'acide sulfurique donne du sulfate d'aluminium :
\[ \ce{Al(OH)3 + H2SO4 -> Al2(SO4)3 + H2O}\qquad\text{(non pondérée).} \]
Pondère l'équation en traitant le motif \ce{SO4} \textbf{en bloc}. \emph{(Astuce : équilibre Al,
puis le groupe \ce{SO4}, puis H et O par l'eau.)}
\end{exo}

\begin{exo}[une précipitation]
On mélange le chlorure d'aluminium et le nitrate d'argent ; il précipite du chlorure d'argent :
\[ \ce{AlCl3 (aq) + AgNO3 (aq) -> AgCl (s) + Al(NO3)3 (aq)}\qquad\text{(non pondérée).} \]
Pondère l'équation (traite \ce{NO3} en bloc) et donne la somme des coefficients.
\end{exo}

\begin{exo}[équilibrer les charges : équation ionique]
L'aluminium métallique réagit avec les ions \ce{H+} d'un acide :
\[ \ce{Al (s) + H+ (aq) -> Al^3+ (aq) + H2 (g)}\qquad\text{(non pondérée).} \]
\textbf{a)} Pondère les \emph{atomes}.\quad
\textbf{b)} Vérifie que la \textbf{charge électrique totale} est la même des deux côtés — c'est la
condition supplémentaire propre aux équations ioniques.
\end{exo}

\end{document}
